% Sumber: authority/upstream/AlJabr-2-9a5803ff77dd3257484cb177f851a73770a59dd3/appendix1.tex, baris 121--283.
\section{Objek kompak dan kategori presentabel}\label{sec:cpt-objects}
Kekompakan dalam teori kategori merupakan penyempurnaan bersama dari banyak kondisi keterhinggaan dalam aljabar, geometri, dan geometri aljabar. Dalam bagian ini kita menetapkan kategori $\mathcal{C}$; kecuali dinyatakan lain, selalu diasumsikan bahwa $\mathcal{C}$ memiliki semua $\varinjlim$ kecil terfilter.

Pertimbangkan kategori kecil terfilter $I$ dan funktor $\alpha: I \to \mathcal{C}$. Untuk $X \in \Obj(\mathcal{C})$, keluarga morfisme kanonik $\iota_i: \alpha(i) \to \varinjlim \alpha$ menginduksi keluarga $(\iota_i)_*: \Hom(X, \alpha(i)) \to \Hom\left(X, \varinjlim \alpha \right)$, dengan $i \in \Obj(I)$, dan selanjutnya menginduksi pemetaan kanonik
\begin{equation}\label{eqn:I-small-gen}
	\varinjlim_{i \in \Obj(I)} \Hom\left( X, \alpha(i) \right) \to \Hom\left( X, \varinjlim \alpha \right).
\end{equation}

\begin{definition}\label{def:kappa-filtered}
	\index{poset!$\kappa$-terfilter ($\kappa$-filtered)}
	Misalkan $(P, \leq)$ adalah poset tak kosong dan $\kappa$ adalah kardinal tak hingga. Jika setiap himpunan bagian $P_0$ yang memenuhi $|P_0| < \kappa$ memiliki batas atas di dalam $P$, maka $(P, \leq)$ disebut $\kappa$-terfilter.
\end{definition}

Jika $\kappa \leq \kappa'$, setiap poset $\kappa'$-terfilter juga
$\kappa$-terfilter. Dengan mengambil $\kappa := \aleph_0 = \omega$, kita
memperoleh kembali poset terfilter yang diperkenalkan dalam
Contoh \ref{eg:filtered-poset}.

\begin{definition}\label{def:cpt-objects}
	\index{objek kompak@objek kompak (compact object)}
	\index{objek kompak!$\kappa$-kompak}
	Jika $X \in \Obj(\mathcal{C})$ sedemikian sehingga \eqref{eqn:I-small-gen} merupakan isomorfisme untuk setiap kategori kecil terfilter $I$ dan setiap $\alpha: I \to \mathcal{C}$, maka $X$ disebut \emph{objek kompak} dari $\mathcal{C}$.
	
	Untuk kardinal kecil $\kappa$ yang diberikan, jika $I$ selanjutnya dibatasi pada poset kecil $\kappa$-terfilter, maka $X$ yang memenuhi syarat terkait disebut \emph{objek $\kappa$-kompak}.
\end{definition}

Dalam terminologi Definisi \ref{def:create-limit}, $X$ merupakan objek kompak jika dan hanya jika $\Hom(X, \cdot): \mathcal{C} \to \cate{Set}$ mempertahankan $\varinjlim$ kecil terfilter. Jika $\kappa \leq \kappa'$, setiap objek $\kappa$-kompak juga $\kappa'$-kompak.

\begin{remark}\label{rem:filtered-poset}
	Sepintas, cara pendefinisian objek kompak dan objek $\aleph_0$-kompak tampak berbeda: yang pertama melibatkan semua kategori kecil terfilter, sedangkan yang kedua hanya memperbolehkan poset kecil terfilter. Namun, terdapat fakta yang tidak sepenuhnya sepele \cite[Teorema 1.5]{AR94}:
	\begin{center}\begin{minipage}{0.7\textwidth}
		Untuk setiap kategori kecil terfilter $I$, terdapat poset kecil terfilter $(\mathcal{P}, \leq)$ beserta funktor kofinal $\mathcal{P} \to I$.
	\end{minipage}\end{center}
	Dari sini dapat disimpulkan bahwa kekompakan ekuivalen dengan kekompakan-$\aleph_0$. Pembuktiannya akan digariskan pada bagian latihan.
\end{remark}

Dengan menggunakan definisi kategori terfilter dan deskripsi konkret $\varinjlim$ terfilter di dalam $\cate{Set}$ (Proposisi \ref{prop:filtered-union}), diperoleh
\begin{itemize}
	\item pemetaan \eqref{eqn:I-small-gen} surjektif $\iff$ setiap $f: X \to \varinjlim \alpha$ memiliki faktorisasi $X \xrightarrow{f_i} \alpha(i) \xrightarrow{\iota_i} \varinjlim \alpha$;\footnote{Catatan penerjemah: sumber mencetak sasaran panah terakhir sebagai $X$. Karena $\iota_i$ adalah morfisme kanonik menuju kolimit, sasaran yang bertipe benar ialah $\varinjlim\alpha$.}
	\item pemetaan \eqref{eqn:I-small-gen} injektif $\iff$ untuk setiap $i \in \Obj(I)$ dan setiap $f, g \in \Hom(X, \alpha(i))$ yang memenuhi $\iota_i f = \iota_i g$, terdapat morfisme $i \to j$ dalam $I$ sedemikian sehingga citra $f$ dan $g$ di dalam $\Hom(X, \alpha(j))$ sama.
\end{itemize}

\begin{example}\label{eg:Set-cpt}
	Untuk kasus $\mathcal{C} = \cate{Set}$, $X$ merupakan objek kompak jika dan hanya jika $X$ merupakan himpunan hingga. Arah ``jika'' langsung diperoleh dari pembahasan di atas dan deskripsi konkret $\varinjlim \alpha$. Untuk arah ``hanya jika'', sajikan $X$ sebagai $\varinjlim$ terfilter dari himpunan-himpunan bagiannya yang hingga, lalu pertimbangkan $\identity_X$.
\end{example}

\begin{example}\label{eg:Mod-cpt}
	Misalkan $R$ adalah gelanggang. Untuk kasus $\mathcal{C} = R\dcate{Mod}$, $X$ merupakan objek kompak jika dan hanya jika $X$ memiliki presentasi hingga; dengan kata lain, terdapat $a, b \in \Z_{\geq 0}$ dan barisan eksak modul-$R$
	\[ R^{\oplus b} \to R^{\oplus a} \xrightarrow{p} X \to 0 .\]
	Untuk $\varinjlim$ terfilter di dalam $R\dcate{Mod}$, lihat penjelasan dalam \cite[Catatan 6.2.3]{Li1}. Arah ``hanya jika'' diserahkan sebagai latihan. Berikut garis besar arah ``jika''. Ambil citra basis standar dari $R^{\oplus a}$ untuk memperoleh keluarga pembangkit $x_1, \ldots, x_a$ dari $X$. Pertama, kita buktikan bahwa \eqref{eqn:I-small-gen} injektif: misalkan $f, g \in \Hom(X, \alpha(i))$ memenuhi $\iota_i f = \iota_i g$. Untuk setiap $x \in \{x_1, \ldots, x_a \}$, terdapat $i \to j_x$ sedemikian sehingga citra $f(x)$ dan $g(x)$ di dalam $\alpha(j_x)$ sama. Sifat kategori terfilter kemudian memberikan morfisme $i \to j$ yang diinginkan, sehingga citra $f$ dan $g$ di dalam $\Hom(X, \alpha(j))$ sama.

	Selanjutnya kita buktikan bahwa \eqref{eqn:I-small-gen} surjektif. Pertimbangkan $f: X \to \varinjlim \alpha$. Untuk setiap $x \in \{x_1, \ldots, x_a \}$, terdapat $j_x$ dan $m_x \in \alpha(j_x)$ sedemikian sehingga $f(x) = \iota_{j_x}(m_x)$; sekali lagi, sifat terfilter memastikan bahwa $j_x$ dapat dipilih secara bebas dari $x$, dan kita notasikan pilihan itu dengan $j$. Ini menentukan morfisme $m: R^{\oplus a} \to \alpha(j)$ sedemikian sehingga $fp = \iota_j m$. Karena $b$ berhingga, dengan mengambil $j \to i$ yang cukup jauh lalu mengganti $j$ dengan $i$, kita juga dapat memastikan bahwa pembatasan $m$ pada citra $R^{\oplus b}$ identik nol; ini memberikan $f_i$ yang dicari.
\end{example}

Untuk objek kompak yang ditentukan oleh berbagai jenis struktur aljabar, \cite[Contoh 1.2; Teorema 3.12]{AR94} memberikan pembahasan yang lebih lengkap. Sebagai contoh, objek kompak dalam kategori grup $\cate{Grp}$ tepat merupakan grup-grup yang memiliki presentasi hingga.

Kelas penting lain dari objek kompak berasal dari pembenaman Yoneda $h_{\mathcal{D}}: \mathcal{D} \to \mathcal{D}^\wedge$ yang diingat kembali dalam \S\ref{sec:Yoneda}.

\begin{proposition}\label{prop:Ind-cpt}
	Misalkan $\mathcal{D}$ adalah sebarang kategori dan $X \in \Obj(\mathcal{D})$. Maka $h_{\mathcal{D}}(X)$ merupakan objek kompak di dalam $\mathcal{D}^\wedge$.
\end{proposition}
\begin{proof}
	Untuk kategori kecil terfilter $I$ dan funktor $\alpha: I \to \mathcal{D}^\wedge$ yang diberikan, Teorema \ref{prop:Yoneda} dan konstruksi titik demi titik dari $\varinjlim$ di dalam $\mathcal{D}^\wedge$ (lihat \cite[Proposisi 2.7.8]{Li1}, yang menotasikannya dengan $\Yinjlim$) memberikan
	\begin{align*}
		\varinjlim_i \Hom_{\mathcal{D}^{\wedge}}(h_{\mathcal{D}}(X), \alpha(i)) & \simeq \varinjlim_i \left(\alpha(i)(X)\right) \\
		& = (\varinjlim \alpha)(X) \\
		& \simeq \Hom_{\mathcal{D}^\wedge}(h_{\mathcal{D}}(X), \varinjlim \alpha).
	\end{align*}
	Verifikasi rutin menunjukkan bahwa ini sama dengan pemetaan kanonik dalam \eqref{eqn:I-small-gen}. Kekompakan pun terbukti.
\end{proof}

Dengan membatasi lebih lanjut kategori terfilter $I$ dalam \eqref{eqn:I-small-gen}, diperoleh berbagai variasi kekompakan. Pertama-tama kita berfokus pada kardinal reguler yang pernah diperkenalkan dalam \cite[Definisi 1.4.10]{Li1}; di sini diperlukan penguraian yang lebih cermat\footnote{Pembaca yang tidak berminat pada teori himpunan dapat melewati pembahasan berikut.}. Mula-mula, misalkan $\alpha > 0$ adalah sebarang ordinal tak hingga. Pertimbangkan ordinal limit $\theta > 0$ dan barisan ordinal yang naik secara ketat $(a_\beta)_{\substack{\beta: \text{ordinal} \\ \beta < \theta}}$. Jika $\sup_{\beta < \theta} a_\beta = \alpha$, maka barisan ini disebut \emph{kofinal} dengan $\alpha$\index{kofinalitas}. Jika $\alpha$ merupakan ordinal limit, definisikan
\[ \mathrm{cf}(\alpha) := \inf\left\{ \theta > 0: \text{ordinal limit}, \; \exists (a_\beta)_{\beta < \theta}\; \text{seperti di atas dan kofinal dengan $\alpha$} \right\}; \]
$\inf$ ini bermakna dan dicapai oleh suatu $\theta$ seperti yang ditampilkan; untuk $\sup$ dan $\inf$ ordinal, lihat \cite[Teorema 1.2.10]{Li1}. Dengan mengambil $a_\beta := \beta$, segera diperoleh $\mathrm{cf}(\alpha) \leq \alpha$.
\index[sym1]{cf@$\mathrm{cf}(\alpha)$}

Perhatikan bahwa kardinal tak hingga, ketika dipandang sebagai ordinal, selalu merupakan ordinal limit; lihat pembahasan setelah \cite[Lema 1.4.6]{Li1}.

\begin{definition}\label{def:regular-cardinal}
	\index{kardinal!reguler (regular)}
	Jika kardinal tak hingga $\kappa$, ketika dipandang sebagai ordinal, memenuhi $\mathrm{cf}(\kappa) = \kappa$, maka ia disebut \emph{kardinal reguler}.
\end{definition}

\begin{proposition}\label{prop:cf-generalities}
	Misalkan $\alpha > 0$ adalah ordinal limit:
	\begin{enumerate}[(i)]
		\item $\mathrm{cf}(\mathrm{cf}(\alpha)) = \mathrm{cf}(\alpha)$;
		\item jika himpunan bagian tak kosong $S \subset \alpha$ memenuhi $\sup S = \alpha$, maka sebagai ordinal berlaku $|S| \geq \mathrm{cf}(\alpha)$;
		\item $\mathrm{cf}(\alpha)$ merupakan kardinal reguler.
	\end{enumerate}
\end{proposition}
\begin{proof}
	Untuk (i), kuncinya adalah membuktikan $\mathrm{cf}(\alpha) \leq \mathrm{cf}(\mathrm{cf}(\alpha))$. Misalkan $(a_\beta)_{\beta < \theta}$ kofinal dengan $\alpha$ dan $(b(\gamma))_{\gamma < \psi}$ kofinal dengan $\theta$. Maka $\left(a_{b(\gamma)}\right)_{\gamma < \psi}$ kofinal dengan $\alpha$. Dari sini terlihat bahwa $\mathrm{cf}(\alpha) \leq \mathrm{cf}(\theta)$. Sekarang ambil $\theta = \mathrm{cf}(\alpha)$.
	
	Untuk (ii), ambil sebarang bijeksi $f: |S| \to S$. Ingat bahwa $|S|$ merupakan ordinal, sedangkan $\alpha \notin S$. Kita menyatakan bahwa terdapat ordinal limit $\theta$ dan pembenaman pelestari urutan $i: \theta \hookrightarrow |S|$ (sehingga $\theta \leq |S|$), sedemikian sehingga barisan ordinal $\left( f(i(\beta)) \right)_{\beta < \theta}$ di dalam $S$ naik secara ketat dan $\sup_{\beta < \theta} f(i(\beta)) = \sup S = \alpha$.
	
	Memang, ini merupakan penerapan prinsip rekursi transfinit. Secara lebih rinci, ambil
	\begin{align*}
		i(0) & := \inf |S| = 0, \\
		i(n+1) & := \inf\left\{ t \in |S|: f(t) > f(i(n)) \right\}, \quad n \in \omega := \{0, 1, 2, \ldots\}, \\
		i(\omega) & := \inf\left\{ t \in |S|: \forall k < \omega, \; f(t) > f(i(k)) \right\},
	\end{align*}
	dan lanjutkan secara transfinit dengan cara ini hingga berhenti pada $\theta$. Pernyataan tersebut terbukti. Dari sini segera diperoleh $|S| \geq \theta \geq \mathrm{cf}(\alpha)$.
	
	Sekarang pertimbangkan (iii). Mengingat (i), cukup ditunjukkan bahwa $\mathrm{cf}(\alpha)$ merupakan kardinal. Misalkan $\beta > 0$ adalah ordinal limit. Maka (ii) menyiratkan $\beta \geq |\beta| \geq \mathrm{cf}(\beta)$; dengan menyubstitusikan $\beta = \mathrm{cf}(\alpha)$ dan menggunakan (i), diperoleh $\mathrm{cf}(\alpha) = |\mathrm{cf}(\alpha)|$, yakni $\mathrm{cf}(\alpha)$ merupakan kardinal.
\end{proof}

\begin{proposition}[Lihat {\cite[Teorema 5.10]{Je03}}]\label{prop:cf-Konig}
	Misalkan $\lambda$ adalah kardinal tak hingga. Maka $\mathrm{cf}(2^\lambda) > \lambda$.
\end{proposition}
\begin{proof}
	Ini merupakan hasil standar dalam teori himpunan. Ambil ordinal limit $0 < \alpha \leq \lambda$ dan barisan ordinal $(a_\beta)_{\beta < \alpha}$ yang memenuhi $a_\beta < 2^\lambda$. Tuliskan $\kappa_\beta := |a_\beta| < 2^\lambda$. Menurut definisi $\sup$ ordinal dan operasi kardinal,
	\[ \left| \sup_{\beta < \alpha} a_\beta \right| = \left| \bigcup_{\beta < \alpha} a_\beta \right| \leq \sum_{\beta < \alpha} \kappa_\beta. \]
	Lema König untuk operasi kardinal (lihat \cite[Latihan 1.6]{Li1}, atau \cite[Teorema 5.10]{Je03}) memberikan
	\begin{align*}
		\sum_{\beta < \alpha} \kappa_\beta & < \prod_{\beta < \alpha} 2^\lambda = (2^\lambda)^{|\alpha|} \\
		& = 2^{\lambda \cdot |\alpha|} \leq 2^{\lambda \cdot \lambda} = 2^\lambda ;
	\end{align*}
	langkah terakhir menggunakan \cite[Teorema 1.4.8]{Li1}. Ini menunjukkan bahwa $\lambda \geq \mathrm{cf}(2^\lambda)$ tidak mungkin terjadi.
\end{proof}

\begin{lemma}\label{prop:enough-regular-cardinal}
	Misalkan $T$ adalah himpunan kecil. Maka terdapat kardinal kecil reguler $\mu$ sedemikian sehingga $\mu > |T|$.
\end{lemma}
\begin{proof}
	Ingat bahwa semesta Grothendieck $\mathcal{U}$ telah ditetapkan. Persoalan ini hanya bergantung pada $|T|$, sehingga tanpa mengurangi keumuman dapat diasumsikan bahwa $T \in \mathcal{U}$ tak hingga. Dengan demikian, himpunan kuasanya $P(T) \in \mathcal{U}$. Tuliskan $\lambda := |T|$. Proposisi \ref{prop:cf-Konig} menyatakan bahwa $\mathrm{cf}(2^\lambda) > \lambda$. Karena $\mathrm{cf}(2^\lambda)$ merupakan kardinal reguler (Proposisi \ref{prop:cf-generalities} (iii)) dan dapat dibenamkan sebagai himpunan bagian dari $2^\lambda = |P(T)|$, maka $\mu := \mathrm{cf}(2^\lambda)$ adalah yang dicari.
\end{proof}

Sekarang kita kembali ke alur utama teori kategori. Konsep berikut berasal dari \cite{GU71}, sedangkan buku teks rujukannya adalah \cite[Chapters 1, 2]{AR94}.

\begin{definition}[P.\ Gabriel, F.\ Ulmer]\label{def:presentable-cat}
	\index{kategori!$\kappa$-aksesibel ($\kappa$-accessible)}
	\index{kategori!$\kappa$-presentabel ($\kappa$-presentable)}
	\index{funktor!$\kappa$-aksesibel ($\kappa$-accessible)}
	Misalkan kategori $\mathcal{C}$ seperti di atas dan $\kappa$ adalah kardinal kecil reguler.
	\begin{enumerate}
		\item Kategori $\mathcal{C}$ yang memenuhi syarat-syarat berikut disebut \emph{kategori $\kappa$-aksesibel}:
		\begin{compactitem}
			\item $\mathcal{C}$ memiliki semua $\varinjlim$ kecil $\kappa$-terfilter (Definisi \ref{def:kappa-filtered});
			\item terdapat himpunan bagian kecil $S \subset \Obj(\mathcal{C})$ yang terdiri dari objek-objek $\kappa$-kompak (Definisi \ref{def:cpt-objects}), sedemikian sehingga setiap $X \in \Obj(\mathcal{C})$ dapat disajikan sebagai $\varinjlim$ kecil $\kappa$-terfilter dari unsur-unsur $S$.
		\end{compactitem}
		\item Funktor antara kategori-kategori $\kappa$-aksesibel disebut funktor $\kappa$-aksesibel jika mempertahankan $\varinjlim$ kecil $\kappa$-terfilter.
		\item Kategori $\kappa$-aksesibel yang kokomplet disebut \emph{kategori $\kappa$-presentabel}\footnote{Gabriel pernah menyebutnya kategori aljabar; literatur \cite{AR94, GU71} menyebutnya kategori yang presentabel secara lokal, sedangkan di sini istilahnya diubah mengikuti \cite[A.1.1]{Lu09}.}.
	\end{enumerate}

	Jika $\kappa'$ cukup besar relatif terhadap $\kappa$, syarat-syarat yang bersesuaian menjadi lebih longgar daripada syarat untuk $\kappa$. Kategori yang $\kappa$-aksesibel (atau $\kappa$-presentabel) untuk suatu kardinal kecil reguler $\kappa$ disebut kategori aksesibel (atau presentabel). Demikian pula, funktor antara kategori-kategori presentabel disebut funktor aksesibel jika funktor itu $\kappa$-aksesibel untuk suatu kardinal kecil reguler $\kappa$.
\end{definition}

Dapat diperiksa bahwa kategori $\cate{Set}$ dan $R\dcate{Mod}$ yang dibahas dalam Contoh \ref{eg:Set-cpt} dan \ref{eg:Mod-cpt} keduanya presentabel; bahkan, keduanya $\aleph_0$-presentabel. Lihat juga latihan.

Untuk hasil penting berikut, di sini kita hanya dapat memberikan pembuktian sebagian.

% Reference: nLab, page on Adjoint Functor Theorem
\begin{theorem}[P.\ Gabriel, F.\ Ulmer]
	Misalkan $\kappa$ adalah kardinal kecil reguler dan $F: \mathcal{C} \to \mathcal{D}$ adalah funktor antara kategori-kategori presentabel. Maka,
	\begin{enumerate}[(i)]
		\item $F$ memiliki adjoin kanan jika dan hanya jika ia mempertahankan $\varinjlim$ kecil;
		\item $F$ memiliki adjoin kiri jika dan hanya jika ia aksesibel dan mempertahankan $\varprojlim$ kecil.
	\end{enumerate}
\end{theorem}
\begin{proof}
	Arah ``hanya jika'' dari pernyataan (i) sudah dikenal dengan baik, sedangkan arah sebaliknya mengikuti dari teorema funktor adjoin khusus \ref{prop:SAFT} dan fakta-fakta berikut. Pertama, Lema \ref{prop:generator-via-limit} menyiratkan bahwa $S$ dalam Definisi \ref{def:presentable-cat} merupakan keluarga pembangkit dari $\mathcal{C}$; kedua, kategori presentabel selalu koberdaya-baik, lihat \cite[Teorema 1.58]{AR94}.
	
	Pernyataan (ii) mengikuti dari teorema funktor adjoin Freyd \ref{prop:GAFT}; lihat \cite[Teorema 1.66]{AR94} untuk rinciannya.
\end{proof}
